% SOURCE_AUTHORITY: sga5_fr_workpass.tex, lines 13163--13204 (LF-normalized unit).
% SOURCE_SHA256: 791F4EFFC5E02832D5D77ED03518C8156D6F07E4C8238B03545DB93D883FBB28
\begin{prooftext}[Demostración de 7.12]
El hecho de que $\RG_{!V}(F)\in\Ob(\Delta)$ resulta inmediatamente de 7.1 aplicado al complejo $i_!(F)$, donde $i:V\to C$; en efecto,
\[
\RG_{!V}(F)=\RG_C(i_!F).
\]
Para probar que $\RG_V(F)\in\Ob(\Delta)$, se considera el triángulo distinguido
\[
\begin{tikzcd}[column sep=huge,row sep=large]
& Q^\bullet \arrow[dl] &\\
Ri_!(F) \arrow[rr] && Ri_*(F) \arrow[ul],
\end{tikzcd}
\]
donde $Q^\bullet$ está concentrado sobre $Z=C-V$; de ahí, aplicando $\RG_C$, el triángulo distinguido
\[
\begin{tikzcd}[column sep=huge,row sep=large]
& \RG_Z(Q^\bullet|Z) \arrow[dl] &\\
\RG_{!V}(F) \arrow[rr] && \RG_V(F) \arrow[ul].
\end{tikzcd}
\tag{7.16}
\]
Como $\RG_{!V}(F)\in\Ob(\Delta)$, resulta que $\RG_V(F)\in\Ob(\Delta)$ si y solo si $\RG_Z(Q^\bullet|Z)\in\Ob(\Delta)$, es decir, $(Ri_*(F))_{\bar x}\in\Ob(\Delta)$ para todo $x\in Z$. Recuérdese que, para calcular $(Ri_*(F))_{\bar x}$, se considera el cuadrado cartesiano
\[
\begin{tikzcd}
\bar V \arrow[r,hook,"j"] \arrow[d,"u"'] & C_x \arrow[d]\\
V \arrow[r,hook] & C
\end{tikzcd}
\]
donde $C_x$ es la localización estricta correspondiente al punto geométrico $x$; se tiene
\[
Ri_*(F)_{\bar x}\simeq \RG_{\bar V}(u^*(F)),
\]
en virtud de 7.14, y se tiene, por tanto, $(Ri_*(F))_{\bar x}\in\Ob(\Delta)$ y
\[
\clK{\Delta}((Ri_*(F))_{\bar x})=0.
\]
Resulta que
\[
\clK{\Delta}(\RG_Z(Q^\bullet|Z))=0
\]
y el triángulo distinguido (7.16) da
\[
\clK{\Delta}(\RG_{!V}(F))=\clK{\Delta}(\RG_V(F)).
\]
